\documentclass[12pt,notheorems,usepdftitle=false,aspectratio=1610]{beamer}

\usetheme{Boadilla}
%\usetheme{Madrid}
%\usecolortheme{seahorse}

\setbeamertemplate{navigation symbols}{}

%\usefonttheme[onlymath]{serif}
\usefonttheme[]{serif}
%\useinnertheme{circles}
\setbeamertemplate{caption}[numbered]

\usepackage[T2A]{fontenc}			
\usepackage[utf8]{inputenc}			
\usepackage[english,russian]{babel}
\usepackage{textcomp}
\usepackage{gensymb}
\usepackage{amsmath,amsthm,thmtools,mathtext,cmap,multirow,tcolorbox} 				
\usepackage{graphicx,wrapfig,subfig,mathtools,gensymb}
\usepackage[font=small,labelfont=bf]{caption}

\DeclarePairedDelimiter{\norm}{\lVert}{\rVert} 

\theoremstyle{definition} % Or \theoremstyle{plain}, \theoremstyle{remark}
\newtheorem{result}{Результат}[section]

\graphicspath{{./figs/10}}

\title[Лекция 10]{Эпиполярная геометрия и фундаментальная матрица}

\author{Александр Танченко}
\institute{}
\date{2025}

\begin{document}

%------------------------------------------------------------
\begin{frame}
    \titlepage
\end{frame}

%------------------------------------------------------------
\begin{frame}
\frametitle{Постановка задачи}
\begin{figure}
    \centering
    \includegraphics[height=0.5\textheight]{fig9.pdf}
    \hspace{2cm}
    \includegraphics[height=0.5\textheight]{fig8.pdf}
\end{figure}
Известны
\begin{itemize}
    \item внутренние параметры камер $K$, $K'$
    \item соответствующие точки на изображениях $\mathbf{x}_i\leftrightarrow\mathbf{x}'_i$
\end{itemize}
Требуется найти
\begin{itemize}
    \item $(R,\textbf{t})$ -- положение одной камеры относительно другой 
\end{itemize}
\end{frame}

%------------------------------------------------------------
\begin{frame}
\frametitle{Эпиполярная геометрия}
\begin{figure}
    \centering
    \includegraphics[height=0.4\textheight]{fig1.pdf}
    \includegraphics[height=0.4\textheight]{fig2.pdf}
\end{figure}
\begin{itemize}
    \item \textbf{baseline} -- отрезок, соединяющий центры камер $\boldsymbol{C}$, $\boldsymbol{C}'$
    \item \textbf{epipole} -- точка пересечения baseline с плоскостью изображения
    \item \textbf{epipolar plane} -- плоскость, содержащая $baseline$
    \item \textbf{epipolar line} -- прямая пересечения эпиполярной плоскости с плоскостью изображения
\end{itemize}
\end{frame}

%------------------------------------------------------------
\begin{frame}
\frametitle{Эпиполярная геометрия}
\begin{figure}
    \centering
    \includegraphics[height=0.4\textheight]{fig3.pdf}
    \hspace{1cm}
    \includegraphics[height=0.4\textheight]{fig4.pdf}
\end{figure}
\begin{itemize}
    \item каждой точке $\mathbf{x}$ на левом изображении соответствует эпиполярная линия $\mathbf{l}'$ на правом
    \item при изменении 3D точки $\mathbf{X}$, все эпиполярные линии на каждом изображении пересекаются в эпиполярных точках $\textbf{e}$, $\textbf{e}'$
\end{itemize}
\end{frame}

%------------------------------------------------------------
\begin{frame}
\frametitle{Эпиполярная геометрия}
\begin{figure}
    \centering
    \includegraphics[height=0.4\textheight]{fig5.pdf}
    \hspace{0.5cm}
    \includegraphics[height=0.35\textheight]{fig6.pdf}
    \hspace{0.cm}
    \includegraphics[height=0.35\textheight]{fig7.pdf}
\end{figure}
\begin{itemize}
    \item точки на изображениях и соответствующие им эпиполярные линии
    \item эпиполярные линии на каждом изображении пересекаются в эпиполярных точках
    \item в примере эпиполярные точки лежат за границами изображений
\end{itemize}
\end{frame}

%------------------------------------------------------------
\begin{frame}
\frametitle{Фундаментальная матрица (fundamental matrix)}
\begin{figure}
    \centering
    \includegraphics[height=0.4\textheight]{fig11.pdf}
\end{figure}
\begin{itemize}
    \item каждой точке $\mathbf{x}$ соответствует эпиполярная линия $\mathbf{l}'$
    \item это отображение в однородных координатах задается \textbf{фундаментальной матрицей} $F\in\mathbf{R}^{3\times3}$
    $$
        \mathbf{l}'=F\cdot\mathbf{x}
    $$ 
    \item точка $\mathbf{x}'$, соответствующая точке $\mathbf{x}$, лежит на эпиполярной линии $\mathbf{l}'$
    $$
        \mathbf{x}\leftrightarrow\mathbf{x}'
        \qquad\Leftrightarrow\qquad
        \mathbf{x}'^{\,T}\cdot\mathbf{l}'=0
        \qquad\Leftrightarrow\qquad
        \mathbf{x}'^{\,T}F\mathbf{x}=0
    $$
\end{itemize}
\end{frame}

%------------------------------------------------------------
\begin{frame}
\frametitle{Фундаментальная матрица}
\begin{figure}
    \centering
    \includegraphics[height=0.35\textheight]{fig11.pdf}
\end{figure}
\begin{itemize}
    \item эпиполярная линия $\mathbf{l}'$ проходит через точки $\mathbf{e}'$, $\mathbf{x}'$
    $$
        \mathbf{l}'=\mathbf{e}'\times\mathbf{x}'=[\mathbf{e}']_\times\mathbf{x}'\,,\quad
        [\mathbf{e}']_\times=\widehat{\mathbf{e}}'=
        \begin{bmatrix*}[r]
            0 & -e_3 & e_2\\
            e_3 & 0 & -e_1\\
            -e_2 & e_1 & 0
        \end{bmatrix*}\,,\quad    
    $$
    \item точки $\mathbf{x}$, $\mathbf{x}'$, $\mathbf{X}$ и матрицы камер $P$, $P'$ удовлетворяют уравнениям
    $$
        P\cdot\mathbf{X}=\mathbf{x}\,,\qquad
        P'\cdot\mathbf{X}=\mathbf{x}'
    $$
\end{itemize}
\end{frame}

%------------------------------------------------------------
\begin{frame}
\frametitle{Фундаментальная матрица}
\begin{figure}
    \centering
    \includegraphics[height=0.35\textheight]{fig11.pdf}
\end{figure}
\begin{itemize}
    \item пусть $P^+=\left(P^TP\right)^{-1}P^T$ -- псевдообратная матрица $P^+P=I$, тогда
    $$
        P\cdot\mathbf{X}=\mathbf{x}
        \quad\Rightarrow\quad
        \mathbf{X}=P^+\cdot\mathbf{x}
        \quad\Rightarrow\quad
        \mathbf{x}'=P'\cdot\mathbf{X}=P'P^+\cdot\mathbf{x}
    $$
    \item формула для фундаментальной матрицы
    $$
        \mathbf{l}'=[\mathbf{e}']_\times\mathbf{x}'=[\mathbf{e}']_\times P'P^+\cdot\mathbf{x}
        \quad\Rightarrow\quad
        F=[\mathbf{e}']_\times P'P^+
    $$
\end{itemize}
\end{frame}

%------------------------------------------------------------
\begin{frame}
\frametitle{Фундаментальная матрица}
\begin{figure}
    \centering
    \includegraphics[height=0.3\textheight]{fig11.pdf}
\end{figure}
\begin{itemize}
    \item в дальнейшем будем считать, что всемирная система координат совпадает
          с системой координат левой камеры
    $$
        P=K\cdot[I\,|\,\mathbf{0}]\,,\qquad
        P'=K'\cdot[R\,|\,\mathbf{t}]
    $$
    \item тогда, псевдообратная матрица
    $$
        P^+=\left(P^TP\right)^{-1}P^T=
        \begin{bmatrix*}[c]
            K^{-1}\\
            \mathbf{0}^T
        \end{bmatrix*}
    $$
    \item эпиполярные точки
    $$
        \mathbf{e}'=P'\cdot\mathbf{C}=
        P'
        \begin{bmatrix*}[c]
            \mathbf{0}\\
            1
        \end{bmatrix*}=
        K'\mathbf{t}\,,\quad
        \mathbf{e}=P\cdot\mathbf{C}'=
        P
        \begin{bmatrix*}[c]
            -R^T\mathbf{t}\\
            1
        \end{bmatrix*}=
        KR^T\mathbf{t}
    $$
\end{itemize}
\end{frame}

%------------------------------------------------------------
\begin{frame}
    \frametitle{Фундаментальная матрица}
    \begin{result}
        Фундаментальная матрица $F$, 
        для любых двух соответствующих точек на изображениях $\mathbf{x}\leftrightarrow\mathbf{x}'$
        удовлетворяет уравнению
        $$
            \mathbf{x}'^{\,T}F\mathbf{x}=0
        $$
        \vspace{1mm}
    \end{result} 
    Выражения для фундаментальной матрицы
    \begin{align*}
        F&=[\mathbf{e}']_\times P'P^+\\
        &=[\mathbf{e}']_\times K'RK^{-1}=
        K'^{-T}RK^T[\mathbf{e}]_\times\\
        &=\left[K'\mathbf{t}\right]_\times K'RK^{-1}=
        K'^{-T}RK^T\left[KR^T\mathbf{t}\right]_\times
    \end{align*}
\end{frame} 

%------------------------------------------------------------
\begin{frame}
    \frametitle{Свойства фундаментальной матрицы}
    \begin{itemize}
        \item для любых соответствующих точек $\mathbf{x}\leftrightarrow\mathbf{x}'$
        $$
            \mathbf{x}'^{\,T}F\mathbf{x}=0
            \qquad\Leftrightarrow\qquad
            \mathbf{x}^{\,T}F^T\mathbf{x}'=0
        $$
        \item $F$ -- для $(P\,,P')\quad\Leftrightarrow\quad$ $F^T$ для $(P'\,,P)$
        \item эпиполярные линии
        $$
            \mathbf{l}'=F\mathbf{x}\,,\qquad
            \mathbf{l}=F^T\mathbf{x}'
        $$
        \item эпиполярные точки $\mathbf{e}$, $\mathbf{e}'$ удовлетворяют уравнениям
        $$
            F\mathbf{e}=\mathbf{0}\,,\qquad
            F^T\mathbf{e}'=\mathbf{0}
        $$
        \item в случае $P=K\cdot[I\,|\,\mathbf{0}]$, $P'=K'\cdot[R\,|\,\mathbf{t}]$, фундаментальная матрица равна
        $$
            F=\left[K'\mathbf{t}\right]_\times K'RK^{-1}
        $$
        \item $\det(F)=0\quad\Rightarrow\quad\mathrm{rank}(F)=2\quad\Rightarrow\quad\mathrm{dot}(F)=7$
    \end{itemize}
\end{frame} 

%------------------------------------------------------------
\begin{frame}
    \frametitle{Существенная матрица (essential matrix)}
    \begin{itemize}
        \item \textbf{нормализованные координаты}
        $$
            \widehat{\mathbf{x}}=K^{-1}\mathbf{x}\,,\qquad
            \widehat{\mathbf{x}}'=K^{'-1}\mathbf{x}'
        $$
        \item матрицы камер в нормализованных координатах
        \begin{align*}
            \mathbf{x}=P\mathbf{X}\,,\quad
            P=K[I\,|\,\mathbf{0}]
            \qquad&\Rightarrow\qquad
            P=[I\,|\,\mathbf{0}]\,,\quad
            \widehat{\mathbf{x}}=P\mathbf{X}\\
            \mathbf{x}'=P'\mathbf{X}\,,\quad
            P'=K'[R\,|\,\mathbf{t}]
            \qquad&\Rightarrow\qquad
            P'=[R\,|\,\mathbf{t}]\,,\quad
            \widehat{\mathbf{x}}'=P'\mathbf{X}
        \end{align*}
        \item фундаментальная матрица в нормализованных координатах называется существенной
        $$
            \mathbf{x}'^{\,T}F\mathbf{x}=0
            \qquad\Leftrightarrow\qquad
            \widehat{\mathbf{x}}'E\widehat{\mathbf{x}}=0\,,\quad
            E=K'^TFK
        $$
        \item в нормализованных координатах $K=K'=I$, $F=E$
        $$
            F=\left[K'\mathbf{t}\right]_\times K'RK^{-1}
            \qquad\Rightarrow\qquad
            E=\left[\mathbf{t}\right]_\times R
        $$
    \end{itemize}
\end{frame}

%------------------------------------------------------------
\begin{frame}
\frametitle{Свойства существенной матрицы}
$$
    E=\left[\mathbf{t}\right]_\times R
$$
\begin{itemize}
    \item $\mathrm{dot}(E)=5$
    \item SVD разложение существенной матрицы
    $$
        \mathbf{R}^{3\times3}\ni E
        \,\,\mbox{-- essential matrix}
        \quad\Leftrightarrow\quad
        E=U
        \begin{bmatrix*}[c]
            1 & 0 & 0\\
            0 & 1 & 0\\
            0 & 0 & 0
        \end{bmatrix*}
        V^T
    $$
\end{itemize}
\end{frame}

%------------------------------------------------------------
\begin{frame}
\frametitle{Относительное положение $(R,\mathbf{t})$ камер из $E$}
В нормализованных координатах
$$
    P=[I\,|\,\mathbf{0}]\,,\qquad
    P'=[R\,|\,\mathbf{t}]
$$
\begin{result}
    Для существенной матрицы $E=U\mathrm{diag(1,1,0)V^T}$ возможны четыре решения
    $$
        R=UWV^T\,\,\mbox{или}\,\,UW^TV^T\,,\qquad
        \mathbf{t}=\pm\mathbf{u}_3
    $$
    где
    $$
        W=
        \begin{bmatrix*}[c]
            0 & -1 & 0\\
            1 & 0 & 0\\
            0 & 0 & 1
        \end{bmatrix*}
        \,,\quad
        \mathbf{u}_3
        \,\,\mbox{-- третий столбец}\,\,
        U
        \quad\Rightarrow\quad
        \norm{\mathbf{t}}=\norm{\pm\mathbf{u}_3}=1
    $$
    \textbf{Вывод}: зная $E$, можно найти относительное положение камер $(R,\mathbf{t})$ 
    с точностью до величины сдвига
\end{result}
\end{frame}

%------------------------------------------------------------
\begin{frame}
\frametitle{Геометрическая интерпретация четырех решений}
\begin{figure}
    \centering
    \includegraphics[height=0.6\textheight]{fig12.pdf}
    \hspace{2cm}
    \includegraphics[height=0.6\textheight]{fig9.pdf}
\end{figure}
Нужно выбрать то решение, для которого одна
3D точка $\mathbf{X}$ будет находиться перед обеими камерами.
\end{frame}

%------------------------------------------------------------
\begin{frame}
\frametitle{Метод нахождения относительного положения камер}
\begin{figure}
    \centering
    \includegraphics[height=0.4\textheight]{fig8.pdf}
    \hspace{2cm}
    \includegraphics[height=0.4\textheight]{fig1.pdf}
\end{figure}
\begin{itemize}
    \item вычислить фундаментальную матрицу $F$ по соответствиям $\mathbf{x}_i\leftrightarrow\mathbf{x}_i'$
    \item по фундаментальной матрице найти существенную $E=K^{'\,T}FK$
    \item из существенной матрицы $E$ найти относительное положение камер $(R,\mathbf{t})$
    с точностью до длины вектора смещения центра камер $\mathbf{t}$
\end{itemize}
\end{frame}

%------------------------------------------------------------
\begin{frame}
\frametitle{DLT вычисление фундаментальной матрицы}
\begin{itemize}
    \item для каждой пары $\mathbf{x}\leftrightarrow\mathbf{x}'$ соответствующих точек на изображениях
    $$
        \mathbf{x}^{'\,T}F\mathbf{x}=0
        \qquad\Leftrightarrow\qquad
        (x',y',1)
        \begin{bmatrix*}[c]
            f_{11} & f_{12} & f_{13}\\
            f_{21} & f_{22} & f_{23}\\
            f_{31} & f_{32} & f_{33}
        \end{bmatrix*}
        \begin{bmatrix*}[c]
            x\\ y\\1
        \end{bmatrix*}
        =0
    $$
    \item для $n\geqslant 8$ соответствий $\mathbf{x}_i\leftrightarrow\mathbf{x}'_i$, $i=1,2,\ldots,n$
    $$
        A\mathbf{f}=
        \begin{bmatrix*}[c]
            x'_1x_1 & x'_1y_1 & x'_1 & y'_1y_1 & y'_1 & x_1 & y_1 & 1\\
            \vdots & \vdots & \vdots & \vdots & \vdots & \vdots & \vdots & \vdots\\
            x'_nx_n & x'_ny_n & x'_n & y'_ny_n & y'_n & x_n & y_n & 1
        \end{bmatrix*}
        \mathbf{f}
        =0
    $$
    где 
    $$
        \mathbf{f}=(f_{11},f_{12},f_{13},f_{21},f_{22},f_{23},f_{31},f_{32},f_{33})^T
    $$
\end{itemize}
\end{frame}

%------------------------------------------------------------
\begin{frame}
\frametitle{DLT вычисление фундаментальной матрицы}
\begin{figure}
    \centering
    \includegraphics[height=0.3\textheight]{fig13.pdf}
\end{figure}
\begin{itemize}
    \item если соответствующие точки $\mathbf{x}_i\leftrightarrow\mathbf{x}'_i$ заданы точно, то $\mathrm{rank}(A)=8$
          и $\mathbf{f}$ находится с точностью до множителя
    \item как правило, координаты точек заданы с погрешностями и в этом случае
    $$
        \mathbf{f}=\mathrm{arg}\min_{\norm{\mathbf{f}}=1}\norm{A\mathbf{f}}
        \quad\Rightarrow\quad
        \mathbf{f}\,\,\mbox{-- последний столбец}\,\, V\quad
        \left(A=UDV^T\right)
    $$
    \item необходимо удовлетворить условию $\mathrm{rank}(F)=8$
    $$
        F=U\mathrm{diag}(r,s,t)V^T
        \qquad\mapsto\qquad
        F=U\mathrm{diag}(r,s,0)V^T
    $$
\end{itemize}
\end{frame}

\end{document}
